\chapter{系统测试}

系统测试属于检验平台实现质量、发现潜在缺陷的重要环节。本章会从功能、性能以及安全这三个维度对平台开展测试，并结合相应的测试结果进行分析。

\section{测试环境与工具}

\subsection{测试环境配置}

系统测试工作是在本地开发环境当中进行的，具体的环境配置情况如表~\ref{tab:5-1}~所示。

\begin{table}[H]
    \centering
    \caption{测试环境配置}
    \label{tab:5-1}
    \begin{tabularx}{\linewidth}{lX}
        \toprule
        项目 & 配置详情 \\
        \midrule
        操作系统 & Windows 11 \\
        Node.js 版本 & v20.x LTS \\
        MySQL 版本 & 8.0 \\
        浏览器 & Google Chrome 122 \\
        后端运行端口 & 3000 \\
        前端运行端口 & 5173 \\
        \bottomrule
    \end{tabularx}
\end{table}

\subsection{测试工具}

\textbf{Postman。}主要用来对后端 API 接口进行功能测试，借助设置不同的请求参数、携带或不携带 Token 的请求头以及请求体，检查接口返回的结果是否与预期相契合。Postman 所提供的环境变量功能可以用于管理 Token 值，在登录接口返回 Token 之后会自动写入环境变量，后续请求也会随之自动携带，由此使测试流程得到简化。

\textbf{浏览器开发者工具。}主要用于对前端交互过程以及网络请求开展测试与跟踪，重点对 SSE 长连接能否成功建立、数据流能否正常传输进行检查；在 Network 面板的 EventStream 选项卡当中，可以逐条查看推送内容，同时还可核对各页面状态更新是否得到正确触发，并检查控制台当中的错误日志。


\section{功能测试}

功能测试主要借助黑盒测试方法来进行，按照预先设定好的测试用例，对各个功能模块的输入以及输出结果逐项开展验证，以检查其是否和需求规格说明保持一致。

\subsection{学生端功能测试}

学生端属于平台当中功能相对最为集中的使用端，因此在测试过程中，需要把注册登录、咨询师浏览、预约完整流程、AI 对话以及公告查阅等核心场景纳入覆盖范围。除了对基础正向流程开展验证之外，还需要重点检查重复邮箱注册、预约时间冲突等边界输入，以及越权访问等异常情形下系统错误处理是否正确。

本小节以学生端的核心业务流程为对象来设计测试用例，覆盖用户注册与登录、咨询师浏览与筛选、预约发起与取消、AI 对话以及公告查阅等主要功能场景，并重点对各接口在合法输入及边界输入条件下的响应表现是否契合需求规格开展验证，测试结果如表~\ref{tab:5-2}~所示。

\begin{table}[H]
    \centering
    \caption{学生端功能测试用例}
    \label{tab:5-2}
    \begin{tabularx}{\linewidth}{clXXXc}
        \toprule
        编号 & 测试功能 & 测试输入 & 预期结果 & 实际结果 & 状态 \\
        \midrule
        S-01 & 用户注册 & 合法用户名、邮箱、密码 & 注册成功，返回 Token & 符合预期 & 通过 \\
        S-02 & 重复邮箱注册 & 已存在的邮箱 & 返回 400，邮箱已注册 & 符合预期 & 通过 \\
        S-03 & 用户登录 & 正确用户名和密码 & 登录成功，返回 JWT & 符合预期 & 通过 \\
        S-04 & 错误密码登录 & 正确用户名+错误密码 & 返回 401，密码错误 & 符合预期 & 通过 \\
        S-05 & 浏览咨询师列表 & 携带有效 student Token & 返回可用咨询师列表 & 符合预期 & 通过 \\
        S-06 & 发起预约 & 合法预约时段、咨询师 ID & 创建 pending 状态预约 & 符合预期 & 通过 \\
        S-07 & 时间冲突预约 & 与已有预约时段重叠 & 返回 400，时间冲突 & 符合预期 & 通过 \\
        S-08 & 取消预约 & 自己的 pending 预约 ID & 状态更新为 cancelled & 符合预期 & 通过 \\
        S-09 & AI 助手对话 & 发送情绪倾诉文本 & SSE 流式返回 AI 回复 & 符合预期 & 通过 \\
        S-10 & 多轮 AI 对话 & 连续发送 3 条相关消息 & AI 回复体现上下文理解 & 符合预期 & 通过 \\
        S-11 & 提交咨询评价 & 已完成会话 ID、评分 4 & 评价写入成功 & 符合预期 & 通过 \\
        S-12 & 查看公告列表 & 携带有效 student Token & 返回公告列表，倒序排列 & 符合预期 & 通过 \\
        \bottomrule
    \end{tabularx}
\end{table}

\subsection{咨询师端功能测试}

\begin{table}[H]
    \centering
    \caption{咨询师端功能测试用例}
    \label{tab:5-3}
    \begin{tabularx}{\linewidth}{clXXXc}
        \toprule
        编号 & 测试功能 & 测试输入 & 预期结果 & 实际结果 & 状态 \\
        \midrule
        C-01 & 更新个人档案 & 填写专业方向、可预约时段 & 档案更新成功 & 符合预期 & 通过 \\
        C-02 & 查看待审核预约 & 携带有效 counselor Token & 返回所有 pending 预约 & 符合预期 & 通过 \\
        C-03 & 确认预约 & 预约 ID，操作=confirmed & 状态更新为 confirmed & 符合预期 & 通过 \\
        C-04 & 拒绝预约 & 预约 ID，附填拒绝原因 & 状态更新，原因写入 & 符合预期 & 通过 \\
        C-05 & 创建会话记录 & confirmed 状态预约 ID & 创建 in\_progress 会话 & 符合预期 & 通过 \\
        C-06 & 添加案例笔记 & 会话 ID、笔记内容 & 笔记写入，关联会话 & 符合预期 & 通过 \\
        C-07 & 笔记隔离性 & 另一咨询师账号查询笔记 & 返回空列表 & 符合预期 & 通过 \\
        C-08 & 结束会话 & 进行中的会话 ID & 状态为 completed，时间写入 & 符合预期 & 通过 \\
        \bottomrule
    \end{tabularx}
\end{table}

\subsection{学校管理端功能测试}

\begin{table}[H]
    \centering
    \caption{学校管理端功能测试用例}
    \label{tab:5-4}
    \begin{tabularx}{\linewidth}{clXXXc}
        \toprule
        编号 & 测试功能 & 测试输入 & 预期结果 & 实际结果 & 状态 \\
        \midrule
        A-01 & 添加咨询师账号 & 用户名、邮箱、初始密码 & 账号创建成功，role=counselor & 符合预期 & 通过 \\
        A-02 & 禁用咨询师账号 & 咨询师用户 ID & is\_active 更新为 0 & 符合预期 & 通过 \\
        A-03 & 查看平台统计 & 携带有效 school Token & 返回注册数、预约量等 & 符合预期 & 通过 \\
        A-04 & 发布公告 & 标题、正文内容 & 公告创建，学生端可见 & 符合预期 & 通过 \\
        A-05 & 删除公告 & 公告 ID & 公告删除，列表不再显示 & 符合预期 & 通过 \\
        \bottomrule
    \end{tabularx}
\end{table}

\subsection{功能测试结论}

25 项测试用例已全部通过验证，系统各端功能模块在输入与输出方面的行为均与需求规格说明保持一致，重复注册、时间冲突以及跨用户数据访问等边界情形也都得到了正确处理。

从整体覆盖范围来看，三端正向流程（注册登录、预约发起与审核、AI 对话、公告查阅及评价提交）均已逐项验证通过；异常输入及权限边界方面，重点检验了重复邮箱注册、预约时段冲突以及 student 角色访问 counselor/school 端接口等典型越权场景，系统均返回了预期的错误码和提示信息。特别值得关注的是，案例笔记隔离性测试（C-07）验证了在不同咨询师账号之间确实存在数据隔离：用账号 B 的 Token 查询账号 A 的笔记，接口返回空列表而非错误，这说明过滤逻辑在 Service 层正确落地。总体而言，功能测试结果表明平台已具备投入试用的基本质量水准。


\section{性能测试}

\subsection{普通接口响应时间测试}

借助 Postman 对主要 API 接口进行多次请求测试，并对响应时间的平均值加以记录与统计，结果如表~\ref{tab:5-5}~所示。

\begin{table}[H]
    \centering
    \caption{主要接口响应时间测试结果}
    \label{tab:5-5}
    \begin{tabularx}{\linewidth}{XXlc}
        \toprule
        接口 & 请求类型 & 响应时间均值 & 是否达标（$\leq$500\,ms） \\
        \midrule
        \path{POST /api/v1/auth/login} & 单表查询 + bcrypt 验证 & 约 180\,ms & 达标 \\
        \path{GET /api/v1/student/counselors} & 多表关联查询 & 约 45\,ms & 达标 \\
        \path{GET /api/v1/student/appointments} & 条件过滤 + 关联查询 & 约 38\,ms & 达标 \\
        \path{POST /api/v1/student/appointments} & 冲突校验 + 写入 & 约 62\,ms & 达标 \\
        \path{GET /api/v1/counselor/sessions} & 多表关联查询 & 约 51\,ms & 达标 \\
        \path{GET /api/v1/school/stats} & 多表聚合统计 & 约 120\,ms & 达标 \\
        \bottomrule
    \end{tabularx}
\end{table}

登录接口响应时间相对较长（约 180\,ms），主要耗时来自 bcrypt 的 hash 比较计算——这是密码安全机制的固有代价，cost factor 为 10 时约需 100\,ms 的 CPU 计算，属于预期内的正常开销，对用户体验无实质影响。

\subsection{AI 流式对话响应测试}

针对 AI 心理助手接口进行 10 次对话测试，重点关注首字节时间（TTFB）和完整回复时间：

\begin{itemize}
    \item \textbf{首字节时间（TTFB）：}平均值约在 1.2 秒，最长也没有超过 2.8 秒，整体都在 3 秒的达标线以内。
    \item \textbf{完整回复时间：}会随着回复内容长短的变化而产生较明显的差别，其中较短回复（50 字以内）一般在 2—3 秒左右，中等长度回复（100—200 字）一般在 4—8 秒左右。
    \item \textbf{流式输出体验：}借助 SSE 的分块推送机制，用户不需要等到完整回复全部生成，就可以看到内容被逐字呈现出来，所以主观上的等待感会明显低于轮询方案。
    \item \textbf{稳定性方面的表现：}在 10 次连续开展的测试过程当中，系统没有出现连接中断或者响应超时的情况，SSE 长连接在完整对话周期内一直保持稳定，服务端的流式写入机制以及客户端的逐字渲染机制也能够持续进行协同运行。
\end{itemize}
\vspace{-0.5\baselineskip}
影响 TTFB 的主要因素，主要还是 DeepSeek API 服务端的调度延迟以及本地网络条件，这两项都属于外部依赖；相比之下，系统自身转发链路带来的延迟基本可以忽略不计。

\subsection{并发能力分析}

受测试环境条件的限制，当前还没有借助大规模压测工具来开展模拟验证。站在系统架构的角度来看，Node.js 的事件循环模型使单进程就可以较高效地去处理大量并发 I/O 等待，这样一来能够避免因为线程阻塞而出现性能明显下降；Sequelize 连接池配置为 \texttt{max: 10}，因此可以支撑最多 10 个并发数据库连接；SSE 长连接是以异步流的形式来实现的，所以不会占用线程资源。综合进行评估后可以看出，在 100 用户并发这一设计目标场景下，系统性能可以满足高校单校部署的使用需求。


\section{安全测试}

\subsection{JWT 过期与无效 Token 测试}

运用 Postman 分别去构造三类异常 Token 的测试场景并发起请求，相关测试结果如表~\ref{tab:5-6}~所示。

\begin{table}[H]
    \centering
    \caption{JWT 异常场景测试结果}
    \label{tab:5-6}
    \begin{tabularx}{\linewidth}{lXllc}
        \toprule
        场景 & Token 状态 & 预期响应 & 实际响应 & 结论 \\
        \midrule
        正常请求 & 有效 Token & 200，返回数据 & 符合预期 & 通过 \\
        篡改 Token & 修改 Payload 后重新编码 & 401，Token 无效 & 符合预期 & 通过 \\
        过期 Token & 使用已过期历史 Token & 401，Token 已过期 & 符合预期 & 通过 \\
        无 Token & 不携带 Authorization 头 & 401，未提供认证 & 符合预期 & 通过 \\
        \bottomrule
    \end{tabularx}
\end{table}

在接收到响应之后，系统会把本地凭据自动清除，并且跳转到登录页面，这说明 Token 失效之后的处理链路是完整的。

\subsection{跨角色访问测试}

选用学生角色的有效 Token，分别去尝试访问咨询师端接口以及管理端接口，同时再借助咨询师角色的 Token 去进一步访问管理端接口，相关的测试结果如表~\ref{tab:5-7}~所示。

\begin{table}[H]
    \centering
    \caption{跨角色访问测试结果}
    \label{tab:5-7}
    \begin{tabularx}{\linewidth}{lXllc}
        \toprule
        测试场景 & 请求接口 & 使用角色 & 预期响应 & 结论 \\
        \midrule
        学生访问咨询师接口 & \path{GET /api/v1/counselor/appointments} & student & 403 & 通过 \\
        学生访问管理端接口 & \path{GET /api/v1/school/stats} & student & 403 & 通过 \\
        咨询师访问管理端 & \path{POST /api/v1/school/counselors} & counselor & 403 & 通过 \\
        \bottomrule
    \end{tabularx}
\end{table}

\texttt{roleGuard} 中间件在三端均有效拦截了跨角色的非法访问，403 响应码与预期一致。

\subsection{SQL 注入防护测试}

在登录接口的用户名字段当中注入常见的 SQL 注入语句，像 \texttt{' OR '1'='1}、\texttt{admin'--} 以及 \texttt{'; DROP TABLE users;--}。

系统在整个流程当中都借助 Sequelize ORM 来进行数据库操作，所有查询参数也都经过了参数化处理，所以注入字符串只会被当作普通的字符串值，既不会对数据库查询结果造成影响，也只会返回正常的"用户名或密码错误"提示，没有出现任何数据泄露或者异常行为。

\subsection{密码安全测试}

把数据库 \texttt{users} 表里的 \texttt{password\_hash} 字段直接查询出来之后，可以确认其中保存的是以 \texttt{\$2b\$} 开头的 bcrypt 哈希字符串，这也就说明明文密码在系统各个环节当中都没有被存储或者传输。再用相同密码对同一用户开展重复注册的验证工作后可以看到，两次得到的哈希值并不一致，这是因为 bcrypt 内置了随机盐，由此也说明盐值机制能够正常起作用\cite{provos1999}。

\section*{本章小结}

本章从功能、性能以及安全三个层面对系统开展了较为全面的测试。功能测试方面，共设计了 25 项黑盒测试用例，把学生端、咨询师端以及管理端的核心功能都包括在内，并且测试通过率达到 100\%；性能测试方面，各项普通接口的响应时间均契合 500\,ms 指标，AI 流式对话的首字节时间控制在 3 秒以内，同时从架构角度对系统在 100 并发场景下的可用性进行了论证；安全测试方面，JWT 签名验证、跨角色权限拦截、SQL 注入防护以及密码哈希存储均已验证有效，这说明系统的安全防护体系较为完整。

